\begin{abstract}
On-device ring microgesture recognition can execute inside an inertial sensor
or on a microcontroller, but these placements have rarely been compared under
the same task and hardware constraints. We implement an in-sensor machine
learning core (MLC), a microcontroller convolutional neural network (CNN), and
a microcontroller hyperdimensional classifier (HDC) on one custom ring. Using
longitudinal session data, four-participant cross-user evaluation, held-out
free-living wear, and hardware measurements, we compare recognition and system
cost. CNN achieves the highest within-user macro-F1 (0.978), whereas MLC gives
the strongest cross-user transfer (0.791) and the lowest measured system power
(6.64\,mW versus 14.03\,mW for HDC and 20.58\,mW for CNN). MLC also produces
the fewest free-living false activations. Raw-sample HDC trails in accuracy,
but a statistical-feature ablation raises its within-user macro-F1 from 0.691
to 0.904. These results show that classifier placement and representation
jointly determine the accuracy, robustness, and resource trade-off.
\end{abstract}

\section{Introduction}

Smart rings are emerging as a platform for subtle, always-available input.
Their form factor places the device within reach of the thumb, enabling
microgestures without raising the arm or looking at a screen
~\cite{oura-galaxy-ring,ring-input-survey,microgesture-elicitation}.
Inertial measurement units (IMUs) are well suited to this setting because they
capture linear and rotational finger motion in a compact, low-power package
~\cite{imu-gesture-recognition}. However, continuously streaming raw motion to
a phone consumes radio energy and exposes an uninterrupted motion trace
~\cite{streaming-ring-system}. Recognition should therefore run locally when
possible.

On-device inertial recognition can execute inside the IMU through a machine
learning core (MLC), on the microcontroller (MCU) as a convolutional neural
network (CNN) quantized for deployment, or on the MCU through hyperdimensional
computing (HDC)
~\cite{st-mlc,tflm,kanerva-hdc,hypercam}. These approaches have been evaluated
on different devices, datasets, and tasks, making their deployment trade-offs
difficult to compare directly.

We implement all three architectures on one custom ring and evaluate them with
the same four microgestures and protocol. We compare within-user session
transfer, exploratory leave-one-participant-out transfer, continuous-stream
behavior, and measured power, flash, and RAM. CNN provides the highest
within-user accuracy, whereas MLC has the lowest system cost, the strongest
cross-user transfer, and fewer free-living false activations. Raw-sample HDC
trails both, while a representation ablation recovers much of its
segmented-window accuracy. The result is a measured trade-off among three deployment
choices rather than a single universally best classifier.

\section{Related Work}

IMU-based rings have been studied for thumb-to-finger microgesture input
~\cite{ring1,ring2,ring3}. This work establishes that finger-worn inertial
sensing can capture subtle gestures, but evaluation commonly emphasizes
segmented recordings and a single classifier. Continuous-stream behavior,
particularly false activations during non-gesture activity, is less frequently
reported~\cite{continuous-gesture}.

On-device inertial classification has followed largely separate paths.
Sensor-integrated MLCs execute feature-based decision trees with little MCU
intervention~\cite{st-mlc}. Quantized neural networks provide expressive
temporal models through embedded inference runtimes~\cite{tflm,mcunet}.
HDC provides integer, prototype-based inference without a neural runtime;
HyperCam demonstrates its efficiency on embedded closed-set image tasks, but
does not address continuous IMU streams or heterogeneous null activity
~\cite{kanerva-hdc,hypercam}. Because these systems use different hardware and
protocols, prior results do not provide a controlled comparison of their
accuracy and system costs on the same ring task.

\begin{figure*}[t]
    \centering
    \includegraphics[width=0.95\textwidth]{manuscript/figures/Main.png}
    \caption{Ring platform, gesture set, and on-device classification paths.
    (a) The ring hardware. (b) The four target microgestures. (c) The MLC
    classifies inside the IMU and interrupts the MCU with its result; CNN and
    HDC process windowed FIFO data on the MCU.}
    \label{fig:wide_image}
\end{figure*}

\section{Methods}

\subsection{Hardware}

The ring integrates a Nordic nRF54L15 MCU (128\,MHz Arm Cortex-M33,
256\,KB RAM) and an STMicroelectronics LSM6DSV16X six-axis IMU on a
$15\times15$\,mm PCB. For MCU inference, samples are read through the IMU's
first-in, first-out (FIFO) buffer. The MLC instead executes decision-tree
inference inside the IMU, allowing the MCU to remain in its interrupt-wake
System ON idle state. The ring is worn on the right index finger, and an
asymmetric enclosure fixes its orientation. Bluetooth Low Energy (BLE) is used
only for collection and verification.

\subsection{Study Design}

\subsubsection{Gesture Set}

We study four thumb-scale microgestures---double side tap, double pinch, pinch
hold, and double flick---plus a null class for non-target activity. Double
pinch and pinch hold share a contact posture and differ mainly in temporal
profile, while double flick has a stronger gyroscope signature.

\subsubsection{Data Collection}

Four participants contributed 23 guided sessions, each containing 15
repetitions of each gesture, together with participant-specific null
recordings. The dataset combines a repeated-session longitudinal subset with
broader single-session cross-user data; session counts therefore differ across
participants. Null recordings covered typing, writing, phone use, walking,
object manipulation, conversational hand movement, and unconstrained wear.
The IMU sampled three-axis acceleration and three-axis angular velocity at
120\,Hz and streamed timestamped samples to an iOS application over BLE.

% Add the approved IRB/exemption identifier and consent statement here.

\subsection{Three Classification Deployments}

All methods operate on 128-sample (1.07\,s), six-channel windows. CNN and HDC
advance by 64 samples (0.53\,s), whereas the MLC advances by one non-overlapping
feature window. MLC and CNN learn null as a fifth class; HDC treats null as a
rejection outcome. Continuous results are normalized by recorded time because
the methods have different decision cadences.

\paragraph{In-sensor MLC}
For each fold, class-balanced training windows are exported to ST MEMS Studio
to train a depth-limited decision tree over sensor-supported statistics,
including means, extrema, variance, energy, peak-to-peak range, and vector
magnitudes. Each exported tree is instantiated in an independent offline
evaluator. For the deployed fold, offline decisions were verified against live
sensor outputs. At runtime the MCU waits in System ON idle, wakes on an MLC
interrupt, and reads the predicted class from the sensor.

\paragraph{MCU CNN}
The CNN contains three one-dimensional convolutional layers with 16, 32, and
64 filters, followed by global average pooling and a five-class output. It is
trained per fold using training-fold standardization, rotation, temporal-shift,
and noise augmentation, with early stopping on validation macro-F1. For
deployment, the selected model is converted by full 8-bit integer (int8)
post-training quantization and executed with TensorFlow Lite for
Microcontrollers. The deployed fold retained its float validation and test
macro-F1 after quantization.

\paragraph{MCU HDC}
HDC uses 2,048-dimensional binary spatter codes. Each sample is quantized into
32 levels using integer percentile bounds fitted on training data. Channel and
level hypervectors are bound into temporal trigrams and bundled into one window
hypervector. Class prototypes are accumulated from training examples and
refined with perceptron-style updates. At inference, the four gesture
prototypes compete by Hamming distance; validation-calibrated maximum-distance
and minimum-margin thresholds map low-confidence windows to null. Bit-exact
tests verify agreement between offline and firmware inference.

\subsection{Evaluation Protocol}

Window accuracy is macro-F1 averaged over the four gestures; a prediction of
null counts as an error on the guided test set. For continuous evaluation, an
activation is emitted when $M$ consecutive windows predict the same non-null
gesture, followed by a 1\,s refractory period. We report $M=1$ and $M=2$.
Activations are matched once to guided gestures from motion onset through one
hop after the reference window; every activation on a free-living null stream
is counted as false.

\section{Evaluation}

We compare recognition under session and user variation, deployed system cost,
and continuous-stream behavior.

\subsection{Within-User Recognition Accuracy}

Session-to-session transfer is evaluated on the 20-session longitudinal subset
using five leave-session-out folds. Each fold uses 14 sessions for training,
two for validation, and four for testing; every session appears in a test fold
exactly once. Model fitting and selection use only the corresponding training
and validation data.

Table~\ref{tab:recognition} reports gesture macro-F1. CNN achieves the highest
mean and lowest variation, followed by MLC. Raw-sample HDC is less accurate and
more variable across sessions.

\begin{table*}[t]
    \centering
    \caption{Gesture macro-F1 for five-fold within-user evaluation and
    four-fold cross-user transfer. LOPO mean gives each participant equal
    weight; reported deviations are sample standard deviations.}
    \label{tab:recognition}
    \small
    \begin{tabular}{lcccccc}
        \toprule
        Method & Within-user & LOPO P1 & LOPO P2 & LOPO P3 & LOPO P4
            & LOPO mean \\
        \midrule
        MLC & $0.945\pm0.036$ & 0.913 & 0.787 & 0.726 & 0.737
            & $0.791\pm0.085$ \\
        CNN (float) & $0.978\pm0.015$ & 0.700 & 0.696 & 0.559 & 0.662
            & $0.654\pm0.066$ \\
        HDC & $0.691\pm0.077$ & 0.272 & 0.327 & 0.233 & 0.172
            & $0.251\pm0.065$ \\
        \bottomrule
    \end{tabular}
\end{table*}

In a representation ablation, replacing raw-sample HDC encoding with a
tree-style statistical feature vector increased macro-F1 to $0.904\pm0.039$.
This variant
was evaluated offline and is excluded from the deployment-cost comparison, but
it shows that input representation accounts for much of HDC's segmented-window
accuracy gap.

\subsection{Cross-User Transfer}

We evaluate transfer to unseen users with four leave-one-participant-out
(LOPO) folds. Each fold excludes one participant completely and balances the
remaining training and validation data by participant and gesture class.
Because session counts differ, Table~\ref{tab:recognition} reports each
held-out participant and an equal-participant mean.

MLC achieves the highest macro-F1 for every held-out participant. In this
limited-data setting, its statistical features and shallow tree transfer more
consistently than the raw-window MCU classifiers. Because participant
balancing also caps training examples per participant, the difference from
within-user performance reflects both user variation and limited multi-user
training data. We therefore interpret LOPO as exploratory rather than as a
population-level estimate.

\subsection{On-Device System Cost}

We measure a sensing baseline and representative-fold MLC, HDC, and int8 CNN
firmware with BLE and logging disabled. Current is measured at 4.2\,V over one
60\,s trace after a 15\,s warm-up; flash and statically allocated RAM come from
linked firmware images. Motionless and scripted activity differ by at most
0.02\,mA.

\begin{table}[t]
    \centering
    \caption{Power, flash, and static RAM of the firmware configurations.}
    \label{tab:system_cost}
    \small
    \begin{tabular}{
        l
        S[table-format=1.2]
        S[table-format=2.2]
        S[table-format=6.0, group-separator={,}]
        S[table-format=5.0, group-separator={,}]
    }
        \toprule
        {Config.} & {Current (mA)} & {Power (mW)} & {Flash (B)} & {RAM (B)} \\
        \midrule
        Baseline & 1.86 &  7.81 &  63412 & 27288 \\
        MLC      & 1.58 &  6.64 &  63440 & 21184 \\
        HDC      & 3.34 & 14.03 &  76228 & 41208 \\
        CNN int8 & 4.90 & 20.58 & 170568 & 75568 \\
        \bottomrule
    \end{tabular}
\end{table}

MLC has the lowest power and MCU footprint. Its current is lower than the
baseline because classification moves into the IMU and the MCU no longer
drains the raw-data FIFO; this is a deployment-path difference, not the
isolated cost of adding a tree. HDC occupies an intermediate cost point, while
CNN requires the most power, flash, and RAM because of its model, inference
runtime, and tensor arena.

\subsection{Continuous-Stream Behavior}

We replay the guided test recordings as continuous streams and evaluate false
activations on 2.66 hours of held-out free-living data from the longitudinal
subset, using frozen models and thresholds. At $M=1$, pooled guided-event
recall is 0.844 for MLC, 0.822 for CNN, and 0.505 for HDC. Mean false
activations per hour across the five fold models are 18.78, 84.15, and 728.13,
respectively.

At $M=2$, false activations fall to 3.38, 43.35, and 362.75 per hour, with
guided recall of 0.217, 0.891, and 0.182. These are method-specific operating
points rather than an equal-latency comparison: confirmation adds one 1.07\,s
MLC window but one 0.53\,s hop for CNN and HDC. CNN's measured recall can rise
at $M=2$ because confirmation suppresses an earlier wrong-class activation,
allowing the one-to-one event matcher to associate a later correct activation;
it does not create additional positive windows. Continuous behavior therefore
reveals a recall--false-activation trade-off not captured by segmented-window
macro-F1 alone.

\section{Discussion}

This study is intended to map design trade-offs rather than identify a single
winner. CNN is the strongest choice when within-user accuracy is primary and
its compute and memory costs are acceptable. MLC is preferable under tighter
energy and memory budgets and, in our limited multi-user data, transfers more
consistently across wearers. Raw-sample HDC is not competitive here, but the
feature ablation shows that representation can matter as much as classifier
paradigm; its inexpensive update structure remains relevant when on-device
adaptation is the main constraint, although adaptation was not evaluated here.

Continuous operation adds a second design axis. Similar segmented accuracy or
guided recall can produce very different false-activation rates, while
confirmation changes delay and recall differently at each method's cadence.
On-device ring classifiers should therefore be selected using accuracy,
resource cost, transfer, and continuous behavior together. Our evidence is
bounded by four participants, unequal session counts, one ring placement, and
four gestures, but it provides a controlled starting point for making that
selection on future ring systems.
